\chapter{Conclusion et Perspectives}

\section{Bilan du projet}

MedSys constitue une réponse complète aux défis de la transformation numérique
hospitalière. L'architecture microservices adoptée permet d'atteindre les objectifs
initiaux tout en offrant une base solide pour les évolutions futures.

\begin{table}[H]
\centering
\renewcommand{\arraystretch}{1.4}
\begin{tabular}{|l|c|l|}
\hline
\rowcolor{tabhead}\color{white}\textbf{Objectif} &
\color{white}\textbf{Statut} &
\color{white}\textbf{Module concerné} \\
\hline
Authentification sécurisée JWT + RBAC & \textcolor{medsysgreen}{\bfseries Réalisé} & ms-auth \\
\rowcolor{roweven}
Dossier médical électronique complet & \textcolor{medsysgreen}{\bfseries Réalisé} & ms-patient-personnel \\
Prise de RDV en ligne (anonyme + connecté) & \textcolor{medsysgreen}{\bfseries Réalisé} & MedicalAppointments \\
\rowcolor{roweven}
Liste d'attente automatique & \textcolor{medsysgreen}{\bfseries Réalisé} & MedicalAppointments \\
Notifications email + temps réel & \textcolor{medsysgreen}{\bfseries Réalisé} & MedicalAppointments \\
\rowcolor{roweven}
Gestion du personnel hospitalier & \textcolor{medsysgreen}{\bfseries Réalisé} & ms-personnel \\
Tableau de bord direction & \textcolor{medsysgreen}{\bfseries Réalisé} & medsys-web \\
\rowcolor{roweven}
Export PDF dossier médical & \textcolor{medsysgreen}{\bfseries Réalisé} & ms-patient-personnel \\
QR Code dossier patient & \textcolor{medsysgreen}{\bfseries Réalisé} & ms-patient-personnel \\
\rowcolor{roweven}
Statistiques RDV direction & \textcolor{medsysgreen}{\bfseries Réalisé} & medsys-web \\
\hline
\end{tabular}
\caption{Récapitulatif des objectifs atteints}
\end{table}

\section{Points forts de l'architecture}

\begin{itemize}[itemsep=6pt]
  \item \textbf{Hétérogénéité technologique maîtrisée} : le projet démontre qu'il
    est possible de faire coexister des technologies Java (Spring Boot) et .NET
    (ASP.NET Core) au sein d'un même système cohérent, chaque service utilisant
    la stack la mieux adaptée à son domaine fonctionnel.

  \item \textbf{Isolation des données} : le principe \textit{database per service}
    garantit l'autonomie de chaque microservice — une migration de schéma ou une
    panne base de données n'affecte qu'un seul service.

  \item \textbf{Scalabilité horizontale} : chaque service peut être répliqué
    indépendamment selon la charge (ex. : multiplier les instances de
    MedicalAppointments lors des pics de prises de RDV).

  \item \textbf{Découplage via RabbitMQ} : les événements asynchrones
    (\code{patient.registered}, \code{appointment.created}) permettent aux services
    de communiquer sans dépendance temporelle, renforçant la résilience.

  \item \textbf{Temps réel avec SignalR} : la notification instantanée des créneaux
    libérés améliore significativement l'expérience utilisateur en évitant les
    rechargements manuels.
\end{itemize}

\section{Limites et points d'amélioration}

\begin{infobox}[Points d'amélioration identifiés]
\begin{itemize}[itemsep=4pt]
  \item \textbf{API Gateway} : en l'état, le frontend Vite proxy joue le rôle
    de routeur. En production, un vrai API Gateway (Kong, Traefik, NGINX) serait
    nécessaire pour la gestion centralisée des auth, rate-limiting, SSL/TLS et
    load balancing.

  \item \textbf{Service Discovery} : l'adressage des services est actuellement
    statique (ports codés en dur). Un registre de services (Consul, Eureka) ou
    un orchestrateur (Kubernetes) permettrait l'auto-découverte et la
    haute disponibilité.

  \item \textbf{Observabilité} : intégrer une stack de monitoring complète
    (Prometheus + Grafana pour les métriques, ELK/Loki pour les logs centralisés,
    Jaeger pour le tracing distribué).

  \item \textbf{Tests automatisés} : augmenter la couverture de tests unitaires
    (JUnit, xUnit) et d'intégration (Testcontainers), et mettre en place un
    pipeline CI/CD (GitHub Actions, GitLab CI).

  \item \textbf{Application mobile} : le frontend React Native (Expo) mentionné
    dans le stack technologique reste à connecter aux microservices.

  \item \textbf{Internationalisation} : préparer le système pour une utilisation
    multilingue (fr/ar/en) en vue d'un déploiement dans des établissements à
    clientèle internationale.
\end{itemize}
\end{infobox}

\section{Perspectives d'évolution}

\subsection{Court terme}

\begin{itemize}[itemsep=4pt]
  \item Déploiement Dockerisé complet avec \code{docker-compose.yml} unifié
    orchestrant tous les services, bases de données et RabbitMQ.
  \item Mise en place d'un reverse proxy NGINX en front de tous les services
    pour remplacer le proxy Vite en production.
  \item Intégration d'un système de paiement en ligne pour les consultations
    privées (Stripe, PayPal).
\end{itemize}

\subsection{Moyen terme}

\begin{itemize}[itemsep=4pt]
  \item Migration vers Kubernetes pour l'orchestration des conteneurs et la
    gestion automatique de la scalabilité.
  \item Implémentation d'un service de téléconsultation via WebRTC intégré
    dans le tableau de bord médecin.
  \item Développement du module de facturation et de gestion des assurances.
  \item Application mobile React Native complète pour les patients (iOS / Android).
\end{itemize}

\subsection{Long terme}

\begin{itemize}[itemsep=4pt]
  \item \textbf{Intelligence Artificielle} : aide à la décision médicale,
    détection d'interactions médicamenteuses, prédiction des pics d'affluence.
  \item \textbf{Interopérabilité HL7/FHIR} : conformité aux standards
    internationaux d'échange de données médicales pour l'intégration avec
    d'autres systèmes hospitaliers.
  \item \textbf{Blockchain} : traçabilité inaltérable des accès au dossier
    médical pour une conformité RGPD renforcée.
\end{itemize}

\section{Conclusion générale}

\begin{tcolorbox}[enhanced, colback=medsyslight, colframe=medsysdark,
  fonttitle=\bfseries, title={MedSys -- Une architecture prête pour la production},
  leftrule=5pt, arc=4pt]
MedSys démontre la viabilité d'une architecture microservices hétérogène pour
un système d'information hospitalier moderne. En combinant la robustesse de
Spring Boot, la performance d'ASP.NET Core et la réactivité de React, le
système répond aux exigences fonctionnelles tout en maintenant une séparation
claire des responsabilités.

\vspace{0.4cm}
La modélisation Merise garantit une structure de données rigoureuse et normalisée,
tandis que les diagrammes UML offrent une vision comportementale et structurelle
complète du système. Ce double niveau de modélisation facilite la maintenance,
l'onboarding de nouveaux développeurs et les évolutions futures.

\vspace{0.4cm}
Avec une sécurité transversale (JWT, RBAC, audit logs, brute-force protection),
des communications asynchrones via RabbitMQ et des notifications temps réel
via SignalR, MedSys pose les fondations d'une plateforme hospitalière numérique
robuste, scalable et prête pour un déploiement en environnement de production.
\end{tcolorbox}

